% Generated by scripts/gen_blueprint.py from blueprint/nodes/02-tools.toml. Do not edit.

\chapter{Tools}
\label{chap:02-tools}

General results used by several of the proofs below. Most of them are in the namespace \texttt{Transcendence}, stated for general use rather than for the proof that first needed them. Two facts about exponential polynomials from the four exponentials development, Lemmas~\ref{lem:expPoly_ne_zero} and~\ref{lem:cauchy_estimate_with_zeros}, are here too, because two of the tools rest on them.

The house $\houseof{\alpha}$ of an algebraic number $\alpha$ is the largest absolute value of its conjugates. The length $L(P)$ of a polynomial $P\in\Z[X][Y]$ is the sum of the absolute values of its integer coefficients, and $M(F)$ is the Mahler measure of $F\in\C[X]$.

\section{Algebraic numbers}

% Prove2Me: Transcendence.liouville_house -- https://prove2.me/theorems/88a49c67-0ede-4536-a0bb-9287b4baff0b
\begin{lemma}[Liouville's inequality, house form]
\label{lem:liouville_house}
\lean{Transcendence.liouville_house}
\leanok
Let $K$ be a number field of degree $d$, let $\sigma\colon K\to\C$ be a field embedding, and let $\alpha\in K$, $\alpha\neq0$. If $c\alpha$ is an algebraic integer for some integer $c\neq0$, then
\[ 1\le|c|^{d}\,|\sigma(\alpha)|\,\houseof{\alpha}^{\,d-1}. \]

\emph{Source:} Standard (Liouville's inequality).
\end{lemma}

\begin{proof}
\leanok
The norm of $c\alpha$ is a non-zero integer. Its absolute value is $|c|^{d}|\sigma(\alpha)|$ times the absolute values of the other $d-1$ conjugates of $\alpha$, each at most $\houseof{\alpha}$.
\end{proof}

% Prove2Me: Transcendence.exists_denom_house_monomial_le -- https://prove2.me/theorems/66481554-9f37-4186-9a65-411f85ad2fae
\begin{lemma}[A common denominator for monomials]
\label{lem:exists_denom_house_monomial_le}
\lean{Transcendence.exists_denom_house_monomial_le}
\leanok
Let $K$ be a number field and $(\theta_i)_{i\in I}$ a finite family in $K$. There are a non-zero algebraic integer $b\in K$ and a real number $H\ge1$ such that for every $e\in\N^{I}$ and every $E\in\N$ with $\sum_ie_i\le E$, the number $b^{E}\prod_i\theta_i^{e_i}$ is an algebraic integer whose house is at most $H^{E}$.

\emph{Source:} Standard.
\end{lemma}

\begin{proof}
\leanok
Take $b$ with every $b\theta_i$ an algebraic integer, and $H=1+\houseof{b}+\sum_i\houseof{b\theta_i}$; the house is submultiplicative.
\end{proof}

% Prove2Me: Transcendence.exists_int_pow_repr -- https://prove2.me/theorems/08874ba9-1330-4e88-94e4-bb411e88f3b6
\begin{lemma}[Scaled powers of an algebraic number]
\label{lem:exists_int_pow_repr}
\lean{Transcendence.exists_int_pow_repr}
\leanok
Let $K$ be a field of characteristic zero and let $\alpha\in K$ be algebraic over $\Q$. There are $\ell,A\in\N$ and an integer $L\neq0$ such that for every $n\in\N$ there are integers $r_0,\dots,r_{\ell-1}$ with $|r_l|\le A^{n}$ and
\[ L^{n}\alpha^{n}=\sum_{l<\ell}r_l\,\alpha^{l}. \]

\emph{Source:} Standard.
\end{lemma}

\begin{proof}
\leanok
Take for $L$ the leading coefficient of an integer polynomial vanishing at $\alpha$. Then $L\alpha$ is an algebraic integer, and reducing its powers modulo a monic integer polynomial keeps the coefficients geometrically bounded.
\end{proof}

\section{Siegel's lemma}

% Prove2Me: Transcendence.siegel_entrywise -- https://prove2.me/theorems/be2108be-d4da-41d0-89c3-0a31d75bb6f8
\begin{lemma}[Siegel's lemma with an entrywise bound]
\label{lem:siegel_entrywise}
\lean{Transcendence.siegel_entrywise}
\leanok
Let $A$ be an $m\times n$ matrix with integer entries, where $n\ge1$ and $2m\le n$, and suppose that every entry satisfies $|A_{ab}|\le B$ for some real $B\ge1$. Then there is $t\in\Z^{n}$, $t\neq0$, with $At=0$ and $|t_b|\le nB$ for every $b$.

\emph{Source:} Siegel's lemma \cite[Ch.~I, \S2]{Lang66}, in the form Mathlib states it for integer matrices.
\end{lemma}

\begin{proof}
\leanok
Mathlib's Siegel lemma bounds a solution by $(n\max(1,\|A\|))^{m/(n-m)}$, and $2m\le n$ makes the exponent at most one.
\end{proof}

\section{Exponential polynomials}

% Prove2Me: FourExp.expPoly_ne_zero -- https://prove2.me/theorems/62105f9b-bb74-4229-93da-4357e32498b8
\begin{lemma}[Distinct frequencies]
\label{lem:expPoly_ne_zero}
\lean{Diaz.expPoly_ne_zero}
\leanok
Let $\omega_1,\dots,\omega_l\in\C$ be pairwise distinct, let $q_1,\dots,q_l\in\N$, and let $b_{j,i}\in\C$ ($1\le j\le l$, $0\le i<q_j$) be not all zero. Then the function
\[ f(z)=\sum_{j=1}^{l}\sum_{i=0}^{q_j-1}b_{j,i}\,z^{i}e^{\omega_jz} \]
takes a non-zero value at some point of $\C$.

\emph{Source:} Classical; see \cite{Tij71} or \cite[\S4]{W71}.
\end{lemma}

\begin{proof}
\leanok
Induction on $l$: multiply by $e^{-\omega_lz}$ and differentiate $q_l$ times. This removes the last block and multiplies the leading coefficient of every other block by $(\omega_j-\omega_l)^{q_l}\neq0$.
\end{proof}

% Prove2Me: Transcendence.expSum_first_nonvanishing -- https://prove2.me/theorems/fed847b5-9f11-481c-86d6-09bccacd99ed
\begin{lemma}[First non-vanishing derivative]
\label{lem:expSum_first_nonvanishing}
\lean{Transcendence.expSum_first_nonvanishing}
\leanok
Let $F(z)=\sum_{i\in I}c_ie^{\rho_iz}$ be a finite exponential sum with pairwise distinct $\rho_i\in\C$ and coefficients $c_i\in\C$ not all zero. Let $m\ge1$ and $n\in\N$, and suppose that $F^{(k)}(j)=0$ for all $j\in\{1,\dots,m\}$ and $k<n$. Then there are $r\ge n$ and $l_0\in\{1,\dots,m\}$ with $F^{(r)}(l_0)\neq0$ and $F^{(k)}(j)=0$ for all $j\in\{1,\dots,m\}$ and $k<r$.

\emph{Source:} Standard.
\end{lemma}

\begin{proof}
\leanok
\uses{lem:expPoly_ne_zero}
By Lemma~\ref{lem:expPoly_ne_zero} the function $F$ is not identically zero, so it vanishes to finite order at each of $1,\dots,m$; take for $r$ the least of these orders.
\end{proof}

% Prove2Me: FourExp.cauchy_estimate_with_zeros -- https://prove2.me/theorems/840a447e-cb2b-4bc8-9bc7-4496079c88d7
\begin{lemma}[Cauchy's estimate with zeros]
\label{lem:cauchy_estimate_with_zeros}
\lean{Diaz.cauchy_estimate_with_zeros}
\leanok
Let $F$ be an entire function, $c\in\C$, $\rho\ge0$ and $R>\rho+1$. Let $S$ be a finite set of points $z$ with $|z-c|\le\rho$, and put $\sigma=\sum_{z\in S}\ord_zF$. If $|F(z)|\le M$ whenever $|z-c|=R$, then for every $s\in\N$
\[ |F^{(s)}(c)|\le s!\,\frac{R}{R-1}\Bigl(\frac{\rho+1}{R-\rho}\Bigr)^{\sigma}M. \]

\emph{Source:} \cite[\S4, (4.3)]{W71}.
\end{lemma}

\begin{proof}
\leanok
Divide $F$ by $\prod_{\beta\in S}(z-\beta)^{\ord_\beta F}$. The quotient is bounded on $|z-c|=R$, hence on $|z-c|\le1$ by the maximum principle; Cauchy's estimate on the unit circle about $c$ concludes.
\end{proof}

% Prove2Me: Transcendence.expPoly_grid_estimate -- https://prove2.me/theorems/707e3c54-a49b-4eac-9d9e-72e16ae0627b
\begin{lemma}[Cauchy's estimate on a grid of zeros]
\label{lem:expPoly_grid_estimate}
\lean{Transcendence.expPoly_grid_estimate}
\leanok
Let $y_1,\dots,y_l\in\C$ be linearly independent over $\Q$, and let $F(z)=\sum_{k\in s}c_kz^{e_k}e^{\omega_kz}$ be a finite sum with $|\omega_k|\le W$ and $e_k\le E$ for all $k$. Let $t_j\le A_j$ ($1\le j\le l$) and $n$ be natural numbers, and suppose that $F^{(i)}\bigl(\sum_jm'_jy_j\bigr)=0$ for all $i<n$ and all $m'\in\N^{l}$ with $m'_j<t_j$ for every $j$. Let $u\ge2$ and $Z\ge\bigl(\sum_jA_j|y_j|+1\bigr)(u+1)$. Then for every $m\in\N^{l}$ with $m_j<A_j$ for every $j$, every $r\in\N$ and $w=\sum_jm_jy_j$,
\[ |F^{(r)}(w)|\le2\,r!\,(u-1)^{-n\prod_jt_j}\Bigl(\sum_{k\in s}|c_k|\Bigr)Z^{E}e^{WZ}. \]

\emph{Source:} The analytic step of the extrapolation in \cite[Lemme 5]{W73} and in the six exponentials theorem \cite[Ch.~II]{Lang66}; it rests on \cite[\S4, (4.3)]{W71}.
\end{lemma}

\begin{proof}
\leanok
\uses{lem:cauchy_estimate_with_zeros}
The $\prod_jt_j$ grid points are distinct because the $y_j$ are linearly independent over $\Q$. Apply Lemma~\ref{lem:cauchy_estimate_with_zeros} on the circle of radius $(\rho+1)u$ about $w$, where $\rho=\sum_jA_j|y_j|$.
\end{proof}

% Prove2Me: Transcendence.expPoly_iteratedDeriv_le -- https://prove2.me/theorems/305d8079-de0f-43d6-9a80-299b0601e6d1
\begin{lemma}[Derivatives at a point]
\label{lem:expPoly_iteratedDeriv_le}
\lean{Transcendence.expPoly_iteratedDeriv_le}
\leanok
Let $f(z)=\sum_{j\in J}P_j(z)e^{w_jz}$, over a finite set $J$, with $|w_j|\le W$ for some $W\ge0$, and each $P_j\in\C[z]$ zero or of degree less than $q_j$; put $N=\sum_jq_j$. Let $c\in\C$ and $D\ge0$. If $|f^{(s)}(c)|\le D$ for every $s<N$, then
\[ |f^{(n)}(c)|\le D\,(W+1)^{n+N}\qquad\text{for every }n\in\N. \]
The $w_j$ need not be distinct.

\emph{Source:} No source is claimed; an elementary induction.
\end{lemma}

\begin{proof}
\leanok
Induction on $N$. If $q_{j_0}\ge1$, then $g=f'-w_{j_0}f$ has the same form with $q_{j_0}$ lowered by one, and $f^{(n+1)}(c)=g^{(n)}(c)+w_{j_0}f^{(n)}(c)$.
\end{proof}

\section{Integer polynomials}

% Prove2Me: Transcendence.length_sum_le -- https://prove2.me/theorems/1aeffaef-2834-435e-9a3c-abb35d93e1d8
\begin{lemma}[Length of a sum]
\label{lem:length_sum_le}
\lean{Transcendence.length_sum_le}
\leanok
For every finite family $(f_i)_{i\in s}$ in $\Z[X][Y]$,
\[ L\Bigl(\sum_{i\in s}f_i\Bigr)\le\sum_{i\in s}L(f_i). \]

\emph{Source:} Standard.
\end{lemma}

\begin{proof}
\leanok
\end{proof}

% Prove2Me: Transcendence.length_mul_le -- https://prove2.me/theorems/7467fd0b-b862-4a30-bce0-a77999e765b6
\begin{lemma}[Length of a product]
\label{lem:length_mul_le}
\lean{Transcendence.length_mul_le}
\leanok
For all $P,Q\in\Z[X][Y]$, $L(PQ)\le L(P)\,L(Q)$.

\emph{Source:} Standard.
\end{lemma}

\begin{proof}
\leanok
\end{proof}

% Prove2Me: Transcendence.exists_modByMonic_length_le -- https://prove2.me/theorems/49c5dda6-cade-42fa-a97e-6981d9634c42
\begin{lemma}[Reduction modulo a polynomial]
\label{lem:exists_modByMonic_length_le}
\lean{Transcendence.exists_modByMonic_length_le}
\leanok
Let $Q\in\Z[X][Y]$. There is $C\in\N$ such that for all $b,e,n\in\N$ and every $P\in\Z[X][Y]$ with $L(P)\le b$, every coefficient of $X$-degree at most $e$, and $Y$-degree at most $n$, the remainder $P\bmod Q$ of the division by $Q$ in $Y$ satisfies
\[ L(P\bmod Q)\le C^{n}b, \]
and every coefficient of $P\bmod Q$ has $X$-degree at most $e+Cn$. (When $Q$ is not monic in $Y$, the remainder is $P$, as in Mathlib's \texttt{modByMonic}.)

\emph{Source:} Standard.
\end{lemma}

\begin{proof}
\leanok
\uses{lem:length_sum_le,lem:length_mul_le}
\end{proof}

% Prove2Me: Transcendence.norm_resultant_le_max_norm_eval -- https://prove2.me/theorems/2a14f6ae-4226-47c6-a9ba-10c45a47e70b
\begin{lemma}[Liouville's inequality through the resultant]
\label{lem:norm_resultant_le_max_norm_eval}
\lean{Transcendence.norm_resultant_le_max_norm_eval}
\leanok
Let $F,G\in\C[X]$ with $M(F)\ge1$ and $M(G)\ge1$. Then for every $\theta\in\C$
\[ |\Res(F,G)|\le2^{\deg F\deg G}M(F)^{\deg G}M(G)^{\deg F}\max\bigl(|F(\theta)|,|G(\theta)|\bigr). \]

\emph{Source:} \cite[Corollary 3.7]{RW97}, in the complex form of their Lemma 3.4; they describe the inequality as well known.
\end{lemma}

\begin{proof}
\leanok
\end{proof}

% Prove2Me: Transcendence.exists_irreducible_factor_cofactor_bound -- https://prove2.me/theorems/5e9f740b-cc7d-407a-a9a2-308e5c4e41ca
\begin{lemma}[The cofactor of an irreducible factor]
\label{lem:exists_irreducible_factor_cofactor_bound}
\lean{Transcendence.exists_irreducible_factor_cofactor_bound}
\leanok
Let $P\in\Z[X]$ have positive degree and let $\theta\in\C$. Then $P=Q^{e}R$ with $Q,R\in\Z[X]$, $Q$ irreducible of positive degree, $e\ge1$, and
\[ 1\le2^{\deg Q\deg R}M(Q)^{\deg R}M(R)^{\deg Q}|R(\theta)|. \]

\emph{Source:} The multiplicative step of Gel'fond's lemma \cite{Gel52}, as in \cite[\S3, Lemme 2]{W71}.
\end{lemma}

\begin{proof}
\leanok
\uses{lem:norm_resultant_le_max_norm_eval}
Take for $Q$ an irreducible factor of positive degree at which $|Q(\theta)|$ is least, and for $e$ its multiplicity. Lemma~\ref{lem:norm_resultant_le_max_norm_eval} bounds every irreducible factor of $R$ from below at $\theta$, and the bound is multiplicative.
\end{proof}

\section{Transcendence degree}

% Prove2Me: Transcendence.exists_monic_integral_model -- https://prove2.me/theorems/f61d5666-eeac-4d3e-b3fd-81eba3d5f7a0
\begin{lemma}[A monic integral model]
\label{lem:exists_monic_integral_model}
\lean{Transcendence.exists_monic_integral_model}
\leanok
Let $\omega\in\C$ be transcendental and let $\theta\in\C$ be algebraic over $\Q(\omega)$. There are $\omega_1\in\C$ and $Q\in\Z[X][Y]$, monic in $Y$ of degree $d\ge1$, with $Q(\omega,\omega_1)=0$ and minimal there: no non-zero $A\in\Z[X][Y]$ with $\deg_YA<d$ vanishes at $(\omega,\omega_1)$. Moreover $\omega_1=v(\omega)\,\theta$ for some $v\in\Z[X]$ with $v(\omega)\neq0$.

\emph{Source:} Standard; the normalisation at the start of \cite[\S III]{W73}.
\end{lemma}

\begin{proof}
\leanok
Clear the denominators of the minimal polynomial of $\theta$ over $\Q(\omega)$ with one $v\in\Z[X]$, and rescale $\theta$.
\end{proof}

% Prove2Me: Transcendence.exists_monic_integral_model_presentation -- https://prove2.me/theorems/9a72eeea-1c1f-42b8-9221-24234f891350
\begin{lemma}[One model for finitely many numbers]
\label{lem:exists_monic_integral_model_presentation}
\lean{Transcendence.exists_monic_integral_model_presentation}
\leanok
Let $\omega\in\C$ be transcendental and let $(z_i)_{i\in I}$ be a finite family of complex numbers algebraic over $\Q(\omega)$. There are $\omega_1\in\C$ and $Q\in\Z[X][Y]$, monic in $Y$ of degree $d\ge1$, with $Q(\omega,\omega_1)=0$ and minimal there in the sense of Lemma~\ref{lem:exists_monic_integral_model}, and $D,E_i\in\Z[X][Y]$ with $D(\omega,\omega_1)\neq0$ and
\[ z_i\,D(\omega,\omega_1)=E_i(\omega,\omega_1)\qquad\text{for every }i. \]

\emph{Source:} Standard; the presentation used in \cite[\S III]{W73}.
\end{lemma}

\begin{proof}
\leanok
\uses{lem:exists_monic_integral_model}
By the primitive element theorem the $z_i$ lie in $\Q(\omega)(\theta)$ for one $\theta$. Apply Lemma~\ref{lem:exists_monic_integral_model} to $\theta$ and clear denominators.
\end{proof}

% Prove2Me: Transcendence.trdeg_adjoin_le_one_of_isAlgebraic_adjoin -- https://prove2.me/theorems/3a193479-d3d2-4e2b-84bf-6398416dcb95
\begin{lemma}[Transcendence degree at most one]
\label{lem:trdeg_adjoin_le_one_of_isAlgebraic_adjoin}
\lean{Transcendence.trdeg_adjoin_le_one_of_isAlgebraic_adjoin}
\leanok
Let $K\subseteq L$ be fields, $x\in L$, and $S\subseteq L$ a set each of whose elements is algebraic over the ring $K[x]$. Then the $K$-algebra $K[S]$ generated by $S$ has transcendence degree at most one over $K$.

\emph{Source:} Standard.
\end{lemma}

\begin{proof}
\leanok
\end{proof}

% Prove2Me: Transcendence.isAlgebraic_adjoin_of_not_algebraicIndependent_pair -- https://prove2.me/theorems/afa7e962-b8de-4d5b-83b2-b7127343a179
\begin{lemma}[A dependent pair]
\label{lem:isAlgebraic_adjoin_of_not_algebraicIndependent_pair}
\lean{Transcendence.isAlgebraic_adjoin_of_not_algebraicIndependent_pair}
\leanok
Let $R$ be a commutative ring, $A$ a commutative $R$-algebra and $t\in A$ transcendental over $R$. If $t$ and $z\in A$ are not algebraically independent over $R$, then $z$ is algebraic over the subalgebra $R[t]$.

\emph{Source:} Standard.
\end{lemma}

\begin{proof}
\leanok
\end{proof}
